% ===========================================================================
% II. Related Work
% ===========================================================================
\section{Related Work}
\label{sec:related}

\subsection{Blockchain-Based Reputation Systems for Supply Chains}

The application of blockchain technology to supply-chain reputation has been
studied extensively.  Carboni~\cite{carboni2015feedback} demonstrated that
decentralized feedback mechanisms can replace centralized rating authorities
without sacrificing data integrity.  Sharples and Domingue~\cite{sharples2016blockchain}
explored the use of blockchain in educational credentialing, establishing
the general principle that tamper-evident ledgers can encode trust signals
across organizational boundaries.  More recent work on AM-specific supply
chains~\cite{almaayn2024unified} implemented an integrated provenance-plus-reputation
system on Hyperledger Fabric v3.1.0, achieving sub-54~ms integrated write
latency and 280--304~TPS concurrent throughput, while defending against 9
scripted attack vectors in a static evaluation.  Our work directly builds
upon that system's threat model and extends the evaluation methodology from
static scripted injection to learned adversarial agents.

\subsection{Adversarial Reputation Manipulation}

Sybil attacks on reputation systems were formalized by Douceur~\cite{douceur2002sybil},
who showed that without a centralized authority, any reputation system is
vulnerable to identity multiplication.  Subsequent work proposed stake-based
economic deterrence~\cite{levien2009attack}, Wilson confidence
intervals~\cite{wilson1927probable} to widen uncertainty around sparse
identities, and gossip-protocol-based cross-verification.  Collusion
detection has been studied via graph-theoretic methods~\cite{christodoulou2019collusion}
and reinforcement-learning-based adaptive agents~\cite{ganesh2019reinforcement}.
Evidence-tampering defenses typically rely on cryptographic commitments or
off-chain audit logs~\cite{ferretti2021blockchain}.  Our work operationalizes
all of these attack classes as discrete actions in a MARL environment,
enabling simultaneous empirical comparison under identical training conditions.

\subsection{Multi-Agent Reinforcement Learning for Security Evaluation}

MARL has been applied to network security in several domains.
Lowe et al.~\cite{lowe2017multi} introduced Multi-Agent DDPG (MADDPG), demonstrating
that centralized training with decentralized execution can stabilize
competitive multi-agent environments.  Foerster et al.~\cite{foerster2018counterfactual}
introduced counterfactual multi-agent policy gradients (COMA) for cooperative
settings.  OpenAI Five~\cite{berner2019dota} and AlphaStar~\cite{vinyals2019grandmaster}
demonstrated that MARL agents can discover sophisticated emergent strategies
through self-play, including coordination and deception.

In the security domain, Pham et al.~\cite{pham2021deep} applied MARL to
autonomous penetration testing, while Hammar and Stadler~\cite{hammar2022learning}
trained RL agents for intrusion prevention in computer networks.  Xu et
al.~\cite{xu2020adversarial} used adversarial MARL to evaluate Byzantine
robustness of federated learning.  Closest to our setting, Ngan et
al.~\cite{ngan2020adversarial} applied RL to reputation-based cooperative
networks, finding that agents with higher reward bonuses for adversarial
actions can undermine group cooperation when defenses are absent.  Our work
differs in explicitly modeling layered defense mechanisms and quantifying
block rates under each defense class, providing actionable security
characterization rather than a proof of vulnerability.

\subsection{Multi-Agent PPO}

Yu et al.~\cite{yu2021surprising} showed that applying PPO independently to
each agent (IPPO) with a shared policy network is surprisingly competitive
with more complex MARL algorithms in cooperative settings.  MAPPO adds
centralized value functions to this framework, providing variance reduction
without sacrificing decentralized execution.  We adopt MAPPO with a shared
actor-critic network (MLP 256$\to$128$\to$ReLU), using GAE advantage
estimation~\cite{schulman2015high} with $\lambda=0.95$, clip ratio $0.2$,
and higher entropy coefficient ($0.05$) to encourage exploration of all 12
actions during early training.

\subsection{PettingZoo and AEC Environments}

Terry et al.~\cite{terry2021pettingzoo} introduced PettingZoo, a standard
API for multi-agent environments using the Agent Environment Cycle (AEC)
formalism~\cite{liu2019aec}.  AEC environments model sequential agent
selection with explicit termination and truncation signals, matching the
turn-based structure of blockchain transaction processing where one agent
submits a transaction at a time.  Our environment follows the AEC contract
strictly, passing \texttt{None} actions for terminated agents and advancing
the agent selector through the full population each episode step.
